\subsection{A whole-output failure already for one-plus-three}

The \fid{F67} obstruction also occurs for the
actual restriction \(\Delta^{(1,3)}\), at the
same total factor count \(k=4\) where
\fid{F66}'s different two-plus-two tensor has
a positive model. The strand rank is still
\(r=3\). This supplement makes no claim that
four is a minimal threshold among all cases.

For two balanced source modules write
\[
 N_1=\begin{pmatrix}m&a\\0&0\end{pmatrix},\quad
 N_2=\begin{pmatrix}0&0\\a&m\end{pmatrix},\qquad
 L_1=\begin{pmatrix}n&b\\0&0\end{pmatrix},\quad
 L_2=\begin{pmatrix}0&0\\b&n\end{pmatrix}.
\]
Take the actual modules
\(\check M^{(1)}\) and \(\check M^{(3)}\), with
\begin{equation}\label{eq:oneThreeActualInputs}
 m=\delta,\quad a=1,\quad
 n=\delta^3,\quad b=3D-4,\qquad
 \delta=v+v^{-1},\quad D=\delta^2.
\end{equation}
These are the source equation~(13), p.~8
and Theorem~3.4, p.~12 of
\cite{BlasiakChebyshev}, since their squared
Chebyshev parameters are \(1\) and \((4-3D)^2\).
The generic irreducibility factors are
\[
        m^2-a^2=D-1,\qquad
        n^2-b^2=(D-1)(D-4)^2.
\]
All four input characters belong to
\(\N[v,v^{-1}]\). Finite graded spaces of
these characters supply effective two-state
kernels and actual quadratic maps with
coefficients \(m\) and \(n\), exactly as
in \fid{F67}. The full-Tate functor of
\fid{F61} supplies their individual pure
weight-zero derived realizations under
\(v=Q^s\), where \(s\) is the extension degree.
Neither construction yet supplies the
categorical restriction on their tensor product.

\paragraph{Direct proof on the whole four-dimensional output.}
The actual inclusion of Definition~2.5 and
equations~(8)--(11), p.~6 of
\cite{BlasiakChebyshev} gives
\[
 R_i=N_i\otimes L_i+(mI-N_i)\otimes(nI-L_i).
\]
With \(h=mn\), these are
\begin{equation}\label{eq:oneThreeFullMatrices}
 R_1=\begin{pmatrix}
 h&mb&an&2ab\\0&0&0&-an\\0&0&0&-mb\\0&0&0&h
 \end{pmatrix},\qquad
 R_2=\begin{pmatrix}
 h&0&0&0\\-mb&0&0&0\\-an&0&0&0\\2ab&an&mb&h
 \end{pmatrix}.
\end{equation}
They satisfy \(R_i^2=hR_i\).
Their product has the full diagonal
\[
 m^2n^2-a^2n^2-m^2b^2+4a^2b^2,\quad
        -a^2n^2,\quad -m^2b^2,\quad m^2n^2.
\]
Consequently
\[
 \Tr(R_1R_2)=2m^2n^2-2a^2n^2-2m^2b^2+4a^2b^2.
\]
Substitution of \eqref{eq:oneThreeActualInputs}
and \(A_4=D^2-8D+8\) gives
\begin{equation}\label{eq:oneThreeNegativeTrace}
\begin{split}
 \tau_{1,3}(v)
 &=2D^4-20D^3+84D^2-128D+64\\
 &=D^2(D-2)^2+A_4^2\\
 &=2(v^8+v^{-8})-4(v^6+v^{-6})
       +20(v^4+v^{-4})+20(v^2+v^{-2})+52.
\end{split}
\end{equation}
The coefficients at \(v^{6}\) and \(v^{-6}\)
are \(-4\). Thus no simultaneous basis makes
both named generators matrices in
\(\N[v,v^{-1}]\): their product and ordinary
trace would then have nonnegative coefficients.
The same tensor/direct-sum argument excludes
effective finite graded kernels preserving
the actual compositions and unmarked closed-word
trace on this entire four-dimensional output.
No output simple projector was used.

The formula agrees with the source decomposition
into the \(j=2\) and \(j=4\) simple constituents,
but the direct product calculation proves it
independently of that decomposition.
In contrast, \fid{F66}'s two-plus-two trace
is \(D^4+A_4^2\).
The different inputs change the other constituents'
contribution; the two conclusions are compatible.

\paragraph{Occurrence, specialization and arithmetic scope.}
The proof of Theorem~3.4, equations~(25)--(26),
p.~12, and generic split semisimplicity in
Corollary~3.5, p.~13 of
\cite{BlasiakChebyshev}, place
\(\check M^{(3)}\) in the third tensor power
of \(\check M^{(1)}=M_{(2,1)}\).
Hence the entire tested module is an actual
generic summand of the fourfold ordinary
Specht tensor and therefore of
\[
 (V_1\otimes V_2\otimes V_3\otimes V_4)^{\otimes3},
 \qquad\dim V_i=2,
\]
whose dimension is \(4096\).
We do not assert a negative coefficient in
the trace of that entire ambient tensor.

The \(\check M^{(3)}\) input is the same regular
two-state lattice as in \fid{F67}.
At \(v=1\), its normalized involutions have
product cube
\(\left(\begin{smallmatrix}7&6\\-6&-5\end{smallmatrix}\right)\),
rather than the identity.
It is therefore not identified with an
ambient specialized direct summand.
The generic occurrence and the exact
matrices above need no such identification.

Under either \(v=Q^{s/2}\) or \(v=Q^s\),
every visible root of
\eqref{eq:oneThreeNegativeTrace} has even
point weight, whereas the Euler multiplicities
have both signs. A pure weight-\(w\) point
complex allows a root of weight \(e\) only
in degree \(e-w\), by BBD~5.1.8,
pp.~126--127 \cite{BBD}.
Its visible signs would all be \((-1)^w\).
Eventual power-sum uniqueness excludes this
uninserted pure-point trace, including hidden
roots and Jordan blocks, because a fixed root
has one permitted pure cohomological degree.
A single overall cohomological parity is also
impossible; Tate or parity shifts do not repair
the mixed signs. Inserted traces are outside
this claim.

A bounded mixed Tate witness does exist:
under \(v=Q^{s/2}\), place the \(136\)
positive-coefficient lines of the displayed
Laurent polynomial in degree zero, and four
copies each of \(E(-3)\) and \(E(3)\) in
degree one, for \(E/\Q_\ell\) finite and
\(\ell\nmid Q\).
Its Euler dimension is \(136-8=128=\tau_{1,3}(1)\).
This scalar witness supplies no effective
operation comparison.
Also \(\tau_{1,3}(v)>0\) for positive real
\(v\), by its sum-of-squares expression;
the obstruction is its Laurent coefficients.

% Separate paragraph for Part III.
\paragraph{The \fid{F67} one-plus-three comparison in Grothendieck's language.}
The named action and source inclusion are
finite-dimensional module data, within
EGA~I, Theorem~1.3.7, p.~86,
Corollary~1.3.8, p.~87 and
Theorem~1.4.1, pp.~90--92 \cite{EGA}.
SGA~3, I, 4.7.2--3 and cote~90, scan~17
supply action language \cite{SGA3,Archive},
not a Hopf structure or an effective iteration.
SGA~6, V, \S\S2--3 and cote~133, scan~18
supply the effective/virtual coefficient cone
\cite{SGA6,Archive}, with no noncommutative
\(\lambda\)-ring claim.
The mixed Tate witness is a perfect point
complex in SGA~6, I, Definition~4.7,
pp.~121--122 and Example~4.8, p.~122
\cite{SGA6}; BBD~5.1.5 and~5.1.8,
pp.~125--127 fix the integral-weight purity
convention \cite{BBD}.
The effective and pure operation comparisons
remain extra hypotheses, not published
canonical-basis axioms.
